\scanpage{27}
\begin{center}-4.13-\end{center}
This is the first step of the construction. Now fix some projection embeddings of $\overline V_1$ and $\overline V_2$, say $\overline V_1\subset\mathbb P_{n_1}$, $\overline V_2\subset\mathbb P_{n_2}$. Embed $\mathbb P_{n_1}$ and $\mathbb P_{n_2}$ as disjoint linear subspaces of $\mathbb P_{n_1+n_2+1}$, and consider the two subvarieties $\overline V_i\subset\mathbb P_{n_1+n_2+1}$. Let $\overline V_4$ be the locus of lines in $\mathbb P_{n_1+n_2+1}$ joining birationally corresponding points of $\overline V_1$ and $\overline V_2$ (cf. Figure 8, where all dimensions have been lowered by one). Now the group $\mathbb G_m$ operates in a natural way on $\mathbb P_{n_1+n_2+1}$ so as to leave $\mathbb P_{n_1}$ and $\mathbb P_{n_2}$ pointwise fixed. This induces an action of $\mathbb G_m$ on $\overline V_4$, leaving the subvarieties $\overline V_1$ and $\overline V_2$ pointwise fixed. Now let $\phi:\overline V_4\longrightarrow V_o$ be the natural projection. Also introduce for convenience a hyperplane $H_o\subset V_o$ through $P_1$ and $P_2$ and left fixed by $\sigma_o$ (there are, in fact, exactly two such hyperplanes). Let $H_4$ be its proper inverse image under $\phi$, i.e., the irreducible subvariety of $\overline V_4$ that dominates $H_o$ and is a component of $\phi^{-1}(H_o)$. Let $I\subset\uline{o}_{\overline V_4}$ be the sheaf of ideals vanishing on $H_4$. Define $\overline V_5$ as the I-transform of $\overline V_4$ (cf. EGA 2, \S8). We wish to describe explicitly the structure of $\overline V_5$. Let $\overline V_{12}$ be the birational joint of $\overline V_1$ and $\overline V_2$, be $v_j$ be the m-adic valuation of the function field $k(V_o)$ centered at $P_j$, and let $\ell_{ij}$ be the center of $v_j$ on $\overline V_i$. Then one verifies readily that $\ell_{ij}$ are lines on $\overline V_i$ and that $\overline V_{12}$ is obtained from $\overline V_i$ by blowing up $\ell_{i1}$ and $\ell_{i2}$. But $\overline V_4$ is essentially a (twisted) product of $(\mathbb P_1-2\text{ points})$ by $\overline V_{12}$, with $\overline V_1$ and $\overline V_2$ at the two ends. $\overline V_5$ differs from $\overline V_4$ only at points of $H_4$ which are singular on $\overline V_4$. In fact, the singularities
\mspage{26}
